% =============================================================================
%  1.2 Promises
% =============================================================================
\section{Промисы (Promises)}

\term{Промис}{Promise} — это объект, представляющий результат асинхронной
операции, который может быть доступен сейчас, в будущем или никогда.
Промисы заменяют цепочки коллбэков и позволяют писать читаемый
асинхронный код.

Промис находится в одном из трёх состояний: \inl{pending},
\inl{fulfilled} или \inl{rejected}.  Переход из \inl{pending}
необратим — это свойство называют \term{неизменяемость}{immutability}
результата.

\subsection{Создание промиса}

\begin{code}{javascript}
const fetchUser = (id) => new Promise((resolve, reject) => {
  const user = db.find(id);
  if (user) {
    resolve(user);
  } else {
    reject(new Error(`User ${id} not found`));
  }
});
\end{code}

\subsection{Комбинаторы}

\subsubsection{Promise.all}

Вероятность успешного завершения $n$ независимых промисов
равна $P = \prod_{i=1}^{n} p_i$, где $p_i$ — вероятность успеха $i$-го.

В блочной форме ожидаемое время завершения \inl{Promise.all}:
\begin{equation}
  T_{\text{all}} = \max_{1 \le i \le n} T_i
\end{equation}

\begin{code}{javascript}
const [user, posts] = await Promise.all([
  fetchUser(1),
  fetchPosts(1),
]);
\end{code}

\note{Используйте \inl{Promise.all}, когда запросы независимы
друг от друга — это позволяет выполнять их параллельно и сократить
общее время ожидания.}

\important{Если хотя бы один промис в \inl{Promise.all} отклоняется,
весь результат будет отклонён. Для случаев, когда нужно получить
результаты всех промисов независимо от ошибок, используйте
\inl{Promise.allSettled}.}
